\documentclass[11pt]{amsart}
\usepackage[T1]{fontenc}
\usepackage{lmodern}
\usepackage{amsmath,amssymb,amsthm,mathtools}
\usepackage{booktabs}
\usepackage{xcolor}
\usepackage[colorlinks=true,linkcolor=blue!50!black,citecolor=blue!50!black,urlcolor=blue!50!black]{hyperref}

\newtheorem{theorem}{Theorem}[section]
\newtheorem{lemma}[theorem]{Lemma}
\newtheorem{proposition}[theorem]{Proposition}
\theoremstyle{definition}
\newtheorem{definition}[theorem]{Definition}
\newtheorem{remark}[theorem]{Remark}

\newcommand{\F}{\mathbb F}
\newcommand{\Mat}{M_3(\F_2)}

\title[The least two-variable identity degree is seven]
{The least degree of a two-variable polynomial identity for
\(M_3(\mathbb F_2)\) is seven}
\author{Zijian Zeng}
\address{Institute of Computer Science and Digital Innovation, UCSI University,
Kuala Lumpur 56000, Malaysia}
\email{zijianzeng@foxmail.com}
\email{1002266693@ucsiuniversity.edu.my}
\date{August 30, 2026}
\subjclass[2020]{16R10, 16Z05, 68W30}
\keywords{polynomial identity, finite field, matrix algebra, exact computation,
rank certificate}
\hypersetup{
  pdftitle={The least degree of a two-variable polynomial identity for M3(F2) is seven},
  pdfauthor={Zijian Zeng}
}

\begin{document}

\begin{abstract}
We determine the least degree of a nonzero polynomial in the two-generator free
algebra \(\F_2\langle X,Y\rangle\) that vanishes on every ordered pair of
matrices in \(M_3(\F_2)\).  The answer is seven.  The upper bound is witnessed
by an explicit 18-word inhomogeneous identity and is checked on all
\(512^2=262144\) ordered pairs.  For the lower bound, evaluations of the 127
words of length at most six on 15 matrix pairs contain a nonsingular
\(127\times127\) minor over \(\F_2\).  Both claims are supplied as a compact
exact certificate and are reconstructed by a dependency-free verifier.
\end{abstract}

\maketitle

\section{Introduction}

The Amitsur--Levitzki theorem gives the standard polynomial of degree \(2n\)
as an identity of \(M_n\), and degree \(2n\) is minimal when the number of free
variables is unrestricted \cite{AmitsurLevitzki1950}.  Requiring exactly two
free variables is a different question: substituting two words into a
multilinear identity can make the resulting free-algebra polynomial zero.  Over
a finite field there is an additional distinction, because an inhomogeneous
polynomial can define the zero function without its homogeneous components
doing so separately.

Genov gave an explicit finite basis for the full ideal of polynomial identities
of \(M_3(K)\) over an arbitrary finite field \(K\)
\cite{Genov1981}.  Later surveys place this result alongside the substantially
different characteristic-zero theory \cite{Drensky2021}.  Those sources do not
tabulate the least degree under the simultaneous restrictions of two free
variables, matrix size three, and base field \(\F_2\).  This precise finite
question was recorded as open in the source research packet
\cite{TheoremDBP2746}; that packet explicitly notes that its status report was
not editorially reviewed.

We resolve the finite question exactly.

\begin{theorem}\label{thm:main}
Let \(d_{3,2}\) be the least degree of a nonzero
\(f\in\F_2\langle X,Y\rangle\) such that
\(f(A,B)=0\) for every \((A,B)\in\Mat^2\).  Then
\[
d_{3,2}=7.
\]
\end{theorem}

The proof is computer-assisted but finite and exact.  Its two decisive objects
are printed below: an explicit identity for the upper bound and a 15-pair rank
witness for the lower bound.  The accompanying verifier reconstructs both from
the serialized data using only integer and bit operations.

\section{Conventions and the evaluation map}

The algebra \(\F_2\langle X,Y\rangle\) is unital.  Its basis is the set of all
words in \(X,Y\), including the empty word \(1\).  For a nonzero polynomial
\(f=\sum_w c_w w\),
\[
\deg f=\max\{|w|:c_w=1\}.
\]
Thus lower-degree and constant terms are allowed.

Let \(W_{\le d}\) be the words of length at most \(d\), ordered first by
length and then lexicographically with \(X<Y\).  Evaluation on all matrix pairs
defines an \(\F_2\)-linear map
\[
E_d:\F_2^{W_{\le d}}\longrightarrow
\F_2^{\Mat^2\times\{0,\ldots,8\}},
\]
where the final coordinate indexes matrix entries in row-major order.  The
desired number is the least \(d\) for which \(E_d\) has a nonzero kernel.
Notice that \(|W_{\le6}|=1+2+2^2+\cdots+2^6=127\).

We encode a matrix \(C=(c_{ij})_{0\le i,j<3}\) as
\begin{equation}\label{eq:encoding}
\iota(C)=\sum_{0\le i,j<3}c_{ij}2^{3i+j}\in\{0,\ldots,511\}.
\end{equation}
This is a bijection.  In particular, exhaustive iteration over the square of
this integer interval includes zero, singular, scalar, commuting, and
noncommuting pairs without any case restriction.

\section{The degree-seven identity}

Define the following polynomial, with every displayed coefficient equal to one
in \(\F_2\):
\begin{equation}\label{eq:identity}
\begin{aligned}
f(X,Y)={}&X^3YXY+X^3Y^2X+XYX^3Y+XY^2X^3\\
&+YX^3YX+YXYX^3+X^4YXY+X^4Y^2X\\
&+X^3YX^2Y+X^3Y^2X^2+X^2YX^3Y+X^2Y^2X^3\\
&+XYX^4Y+XY^2X^4+YX^4YX+YX^3YX^2\\
&+YX^2YX^3+YXYX^4.
\end{aligned}
\end{equation}
The 18 words are distinct, and twelve have length seven.  Hence \(f\ne0\) in
the free algebra and \(\deg f=7\).

\begin{proposition}\label{prop:upper}
The polynomial \(f\) in \eqref{eq:identity} satisfies
\(f(A,B)=0\) for every \((A,B)\in\Mat^2\).
\end{proposition}

\begin{proof}
The verifier constructs multiplication from \eqref{eq:encoding}.  Entry
\((i,j)\) of a product is the parity of the bitwise intersection of row \(i\)
of the first matrix with column \(j\) of the second.  For each of the 262144
encoded pairs, it evaluates the prefix closure of the 18 words by induction on
word length and XORs their values.  Every resulting nine-bit integer is zero.

The computation is exhaustive exact arithmetic.  As an independent arithmetic
check, the discovery program constructs multiplication by the three-term scalar
formula for every entry; the two implementations agree on all \(512^2\) matrix
products.  The verifier's compact success record reports 262144 checked pairs
and zero violations.
\end{proof}

\section{The degree-six rank certificate}

For the lower bound, the following are the 15 encoded ordered pairs used in the
certificate.  The encoding is \eqref{eq:encoding}.

\begin{center}
\begin{tabular}{c@{\qquad}rr@{\qquad}c@{\qquad}rr@{\qquad}c@{\qquad}rr}
\toprule
\(k\)&\(A_k\)&\(B_k\)&\(k\)&\(A_k\)&\(B_k\)&\(k\)&\(A_k\)&\(B_k\)\\
\midrule
0&198&14&5&282&169&10&508&34\\
1&491&43&6&356&227&11&70&92\\
2&278&478&7&430&285&12&366&324\\
3&208&111&8&211&314&13&440&382\\
4&501&140&9&140&459&14&272&72\\
\bottomrule
\end{tabular}
\end{center}

Let \(R\) be the \(135\times127\) binary matrix obtained by evaluating the
ordered words \(W_{\le6}\) on these pairs, with rows indexed by \((k,e)\),
where \(0\le k<15\) and \(0\le e<9\) is the row-major entry index.  Let \(Q\)
be the square submatrix consisting of all nine rows for \(0\le k<14\) and the
single row \((14,3)\).

\begin{proposition}\label{prop:lower}
Over \(\F_2\), the matrix \(Q\) has rank 127.  Consequently, \(E_6\) is
injective.
\end{proposition}

\begin{proof}
Exact Gaussian elimination gives \(\operatorname{rank}_{\F_2}Q=127\), or
equivalently \(\det Q=1\).  The verifier reconstructs every entry of \(Q\)
from the 15 pairs and the word convention; the square matrix itself is not
trusted as input.

Now suppose that a polynomial of degree at most six vanished on every ordered
matrix pair.  Its 127-dimensional coefficient vector would lie in the kernel of
the restriction represented by \(Q\).  Since \(Q\) is nonsingular, this vector
must be zero.  Thus \(E_6\) is injective.  The nested word bases then also give
injectivity in every smaller degree.
\end{proof}

The number of witness pairs is minimal for a certificate of this form: one pair
provides only nine scalar rows, while \(\lceil127/9\rceil=15\).  For reference,
the recomputed ranks on the same pairs are
\[
\begin{array}{c|rrrrrrr}
d&0&1&2&3&4&5&6\\ \hline
\operatorname{rank}E_d|_{\{(A_k,B_k)\}}&1&3&7&15&31&63&127.
\end{array}
\]

\begin{proof}[Proof of Theorem~\ref{thm:main}]
Proposition~\ref{prop:upper} gives \(d_{3,2}\le7\).
Proposition~\ref{prop:lower} gives \(d_{3,2}>6\).  Hence \(d_{3,2}=7\).
\end{proof}

\section{Reproducibility and scope}

The machine-readable certificate contains the literal support of
\eqref{eq:identity}, the 15 encoded pairs, the 127 selected coordinates, and the
claimed lower-degree ranks.  From the artifact root, the complete check is
invoked by
\begin{verbatim}
python3 src/verify_certificate.py \
  certificates/min_degree_7.json
\end{verbatim}
The verifier uses only the Python standard library.  It validates the schema and
ranges, rejects duplicate or malformed data, reconstructs matrix multiplication
and word evaluation, checks all ordered pairs, and independently rebuilds the
rank minor.  A separate test suite compares both multiplication implementations
on their complete domain and tests altered certificates.  No floating-point
arithmetic, random seed, timeout, or unproved search cutoff enters the proof.

The result is specific to the unital two-generator free algebra over
\(\F_2\).  It does not classify the 48-dimensional sampled space of degree-seven
relations, give a new basis for the full identity ideal, or extend the value to
other matrix sizes or fields.  A search through the cited primary and survey
literature, repeated after obtaining the exact polynomial, did not locate this
precise value; this is a bounded literature-search statement rather than a claim
of absolute priority.

\section*{Computational and AI disclosure}

Exact discovery used independently implemented GF(2) evaluators and elimination
routines.  Generative AI assisted with exploratory code, literature queries,
and manuscript drafting.  The mathematical claims are tied to the explicit
certificate and dependency-free verifier described above.  No proof assistant
was used.

\section*{Data and code availability}

The submission archive includes the certificate, independent verifier, test
suite, and instructions needed to reproduce the result from the serialized
exact data.

\bibliographystyle{amsplain}
\bibliography{references}

\end{document}
